% SOICT 2026 submission — Springer CCIS format (single-blind: keep author names).
% Requires the Springer LNCS/CCIS class `llncs.cls` + `splncs04.bst` from
% https://soict.org/wp-content/uploads/2025/09/Latex-Template-for-Springer.zip
% Max 12 pages excluding references. Figure: convert docs/paper/diagrams/soict_harlstmgat.svg -> PDF.
\documentclass[runningheads]{llncs}
\usepackage{graphicx}
\graphicspath{{diagrams/}{./}}   % figure lives in docs/paper/diagrams/
\usepackage{booktabs}
\usepackage{amsmath}
\usepackage[T1]{fontenc}

\begin{document}

\title{HAR versus Deep and Graph-Attention Models for Multi-Horizon Volatility
Forecasting: A Panel-Correct Study on Vietnamese and U.S. Equities}

\author{Author One\inst{1} \and Author Two\inst{1}}
\authorrunning{Author One et al.}
\institute{Affiliation, City, Country \\ \email{\{author1,author2\}@example.edu}}

\maketitle

\begin{abstract}
We study whether deep sequence models or a cross-sectional graph-attention branch can improve daily
Parkinson-variance forecasting over the Heterogeneous Autoregressive (HAR) baseline, and we evaluate them
with panel-correct inference. Three equity panels are used: VN30 (33 tickers) and VN100 (104 tickers) as the
Vietnamese case study, and a long-history S\&P~500 subset (457 tickers) as a data-rich cross-market check.
Against HAR and GARCH(1,1) we benchmark a pooled LSTM over the three HAR features, an LSTM with a Graph
Attention Network branch over a graphical-lasso partial-correlation graph, a frozen-expert convex HAR--deep
combination, and HAR-anchored residual models, at horizons $\{1,5,10,22\}$ with five seeds. Significance is
assessed with the date-clustered Diebold--Mariano test (one loss differential per calendar date); a naive
per-observation test on a cross-sectionally dependent panel overstates significance by a factor of order
$\sqrt{N}$. Three findings follow. First, HAR is very hard to beat on the small Vietnamese panels: under
date-clustered inference no model, including the forecast combination, significantly beats HAR at any horizon
(closest $p=0.08$). Second, on the large S\&P~500 the deep temporal LSTM (no graph) significantly beats HAR,
by up to $+7.1\%$ QLIKE at $h22$ ($p<0.001$), and the convex combination beats HAR at all four horizons
($p<0.001$); competitiveness grows with panel size. Third, the cross-sectional graph adds no out-of-sample
value on any panel: removing it never improves on the no-graph model, a leakage-safe linear neighbour-signal
regressor adds essentially zero incremental out-of-sample $R^2$, and diagnostics rule out a bug or
overfitting. We report MSE, RMSE, MAE, QLIKE and $R^2$ together, since the ranking is metric- and
panel-dependent.
\keywords{Volatility forecasting \and HAR \and LSTM \and Graph Attention Network \and Diebold--Mariano.}
\end{abstract}

\section{Introduction}
Daily volatility forecasting underpins risk management, position sizing and derivative pricing. The HAR
model of Corsi~\cite{corsi2009} is a strong, parsimonious benchmark: three OLS coefficients on daily,
weekly and monthly aggregates capture much of the long-memory structure. Deep sequence models and graph
neural networks are frequently proposed to capture nonlinear temporal dynamics and cross-sectional
spillovers, but the evidence that they beat HAR out of sample is mixed and positive claims are often
confounded by data-design choices.

We ask two questions, with the Vietnamese equity market as the primary case study and a long-history U.S.
S\&P~500 subset as a data-rich cross-market check: (1) Can deep or HAR-anchored models --- a pooled LSTM over
the three HAR features, a convex HAR--deep combination, or HAR-residual correctors --- beat HAR and GARCH out
of sample under panel-correct inference, and at which horizons? (2) Does adding a cross-sectional Graph
Attention Network branch, with edges estimated by graphical lasso, contribute any out-of-sample value beyond
the same-capacity no-graph model? Our contributions are: (i) a panel-correct multi-horizon benchmark of HAR
against an LSTM, an LSTM-GAT, a convex combination and HAR-anchored residual models on three panels, with
significance by the date-clustered Diebold--Mariano test (a naive per-observation test overstates
significance by a factor of order $\sqrt{N}$ on this cross-sectionally dependent panel); (ii) the main
finding that HAR is not significantly beaten on the small Vietnamese panels, whereas on the data-rich
S\&P~500 the deep temporal LSTM significantly beats HAR (up to $+7.1\%$ QLIKE at $h22$) and the convex
combination beats it at every horizon, so deep-versus-HAR competitiveness scales with data; (iii) a
model-free demonstration that the cross-sectional graph adds no out-of-sample value on any panel, a genuine
null rather than a bug or overfitting; and (iv) per-metric fairness --- reporting MSE, RMSE, MAE, QLIKE and
$R^2$ together and showing the ranking is metric- and panel-dependent.

\section{Related Work}
\textbf{HAR and classical models.} HAR~\cite{corsi2009} approximates the long memory of realized volatility
with a daily/weekly/monthly cascade; recent work reports that rolling-window and specification choices,
rather than model class, often drive apparent gains over HAR~\cite{audrino2025}. HARQ~\cite{harq2016}
exploits measurement error, and GARCH-family models remain standard conditional-variance benchmarks. We use
HAR and GARCH(1,1) as the baselines to beat.
\textbf{Deep and graph models.} LSTMs~\cite{lstm1997} are widely applied to financial series, and graph
neural networks model cross-asset spillovers, e.g.\ Graph Attention Networks~\cite{gat2018}, multivariate
realized-volatility GNNs~\cite{chen2022} and spillover-aware GNN-HAR hybrids~\cite{gnnhar2025}; evidence is
mixed. \textbf{Graph construction.} Rather than a market-factor-dominated correlation graph, we estimate a
sparse partial-correlation graph with graphical lasso~\cite{glasso2008} on training rows only.
\textbf{Evaluation.} Volatility is latent, so we use the proxy-robust QLIKE loss~\cite{patton2011} alongside
MSE, with significance by the Diebold--Mariano test~\cite{dm1995} and the Harvey--Leybourne--Newbold
small-sample correction~\cite{hln1997}.

\section{Method}
\textbf{Target and features.} The target is the Parkinson variance~\cite{parkinson1980} at day $t{+}h$ for
horizons $h\in\{1,5,10,22\}$ (the primary date-clustered study; the descriptive per-observation and snapshot
studies use $h\in\{1,5\}$), a non-negative point forecast. Following HAR, three features per node are used: the daily
Parkinson variance at $t$, its 5-day rolling mean (weekly) and its 22-day rolling mean (monthly). The same
features drive every model, so differences reflect functional form, not the information set.

\textbf{Main model: per-observation LSTM.} The main model (``LSTM'') is a two-layer LSTM (hidden 64) over the
lookback window of the three features; a single pooled model is trained over all tickers, each standardized
by its own scaler fit on training rows only, with a linear output inverse-transformed for evaluation.
Overfitting is controlled by dropout, weight decay, gradient clipping, a ReduceLROnPlateau schedule and early
stopping. Training examples are formed by \emph{per-observation pooling}: every (ticker, window) pair is an
independent example and the chronological 80/10/10 split is applied within each ticker's own history,
yielding tens of thousands of windows and interleaving each ticker's regimes.

\textbf{Graph variant: LSTM+GAT (ablation).} To test whether a cross-sectional graph helps, a Graph Attention
Network branch~\cite{gat2018} is added (Fig.~\ref{fig:arch}): a per-node LSTM encodes the temporal window,
and a GAT reads the node features at $t$ and attends over a fixed Top-5 graphical-lasso partial-correlation
graph estimated on training rows only; the branches are concatenated into an MLP head. The leave-one-out
``LSTM (w/o GAT)'' removes the spatial branch. A GAT operates over a common-date cross-section, so LSTM+GAT
requires a \emph{common-date snapshot} design (only dates on which all nodes are present) with a single
global 80/10/10 split; we thus report the graph ablation as a separate study to avoid confounding.

\textbf{Baselines.} HAR is a pooled OLS on the three features refit on the same training rows; GARCH(1,1) is
a per-ticker conditional-variance model.

\textbf{Protocol.} Each configuration runs five seeds $\{42,123,2026,7,2024\}$, up to 20 epochs with early
stopping. We report MSE, RMSE, MAE, QLIKE and $R^2$, seed-averaged, with a shared QLIKE floor $10^{-8}$; the
primary decision metric is QLIKE. Significance is assessed with the \emph{date-clustered} Diebold--Mariano
test: the per-observation loss differential is collapsed to one value per calendar date before the statistic
(HLN correction, HAC lag $h{-}1$) is computed. This is the panel-correct treatment for a cross-section that
shares each trading date; a naive per-observation test treats the effective sample as the number of tickers
times the number of dates and inflates the statistic by a factor of order $\sqrt{N}$ (roughly six on VN30,
ten on VN100). All significance statements use the date-clustered test unless a result is explicitly labelled
per-observation.

\section{Data}
Three panels are used with the Parkinson target and three HAR features: \textbf{VN30} (33 tickers, primary,
Vietnam), \textbf{VN100} (104, vnstock), and \textbf{S\&P~500} ($\sim$500, U.S. robustness; Yahoo-sourced,
only aggregate metrics stored). The main study uses per-ticker 80/10/10 splits and a pooled model; the graph
study uses common-date fixed-node snapshots with a global split. The S\&P~500 list is current membership, so
U.S. levels carry survivorship bias; we claim the short-versus-long ordering transfers, not absolute levels.

\section{Experiments}
We run five studies on the same features across VN30, VN100 and the S\&P~500 subset: (1) the
\emph{primary} HAR-anchored ladder E0--E10 under date-clustered inference at $h\in\{1,5,10,22\}$
(Sec.~\ref{sec:haranchor}), the source of the headline results; (2) a \emph{design comparison}, LSTM vs HAR
and GARCH on the per-observation design (lookback 10; $h\in\{1,5\}$), reporting point estimates and naive
per-observation statistics; (3) a \emph{graph-check} ablation, LSTM+GAT vs LSTM (w/o GAT) on the snapshot
design; (4) \emph{lookback} 10 vs 22 (snapshot, VN30); and (5) a \emph{model-free graph screen} of whether
any cross-stock neighbour signal adds out-of-sample value beyond HAR.

\section{Results}
Parkinson variance is of order $10^{-7}$; MSE, RMSE, MAE are scaled (MSE $\times10^{-7}$, RMSE/MAE
$\times10^{-4}$). Row order is HAR $\to$ GARCH $\to$ LSTM (and, for the graph study, adding LSTM (w/o GAT)
$\to$ LSTM+GAT). All values are five-seed test means.

\subsection{Main study --- per-observation LSTM (descriptive point estimates)}
Table~\ref{tab:main} reports QLIKE point estimates and, for the design comparison only, the \emph{naive
per-observation} Diebold--Mariano statistics. On the Vietnamese panels the LSTM's QLIKE point estimates edge
below HAR at the short horizons; the per-observation test flags several as significant, but that test
overstates significance on a cross-sectionally dependent panel, and under the panel-correct date-clustered
test (Sec.~\ref{sec:haranchor}) none of these Vietnamese gaps is significant. Every learned model beats GARCH
by large margins. On the S\&P~500 the LSTM attains both the lower QLIKE and lower MSE at both horizons, and
the date-clustered study (Sec.~\ref{sec:haranchor}) confirms the S\&P~500 beat-HAR result is statistically
significant there. The panel-correct verdicts are in Sec.~\ref{sec:haranchor}.

\begin{table}[t]
\centering
\caption{Main study (per-observation, per-stock 80/10/10). QLIKE (lower is better; 5-seed mean) and
\emph{naive per-observation} Diebold--Mariano LSTM-vs-HAR on QLIKE and MSE, shown for the design comparison
only ($*$: $p<0.05$ on the per-observation test, which overstates significance on this cross-sectionally
dependent panel --- panel-correct verdicts in Table~\ref{tab:haranchor-vn30}, \ref{tab:haranchor-vn100} and
Sec.~\ref{sec:haranchor}; negative favors LSTM).}
\label{tab:main}
\begin{tabular}{llccc c cc}
\toprule
Panel & $h$ & HAR & GARCH & LSTM & & DM QLIKE & DM MSE \\
\midrule
VN30   & 1 & 0.4675 & 0.7434 & \textbf{0.4578} & & $-3.60^{*}$ & $-0.69$ \\
VN30   & 5 & 0.5514 & 0.7409 & \textbf{0.5458} & & $-0.98$      & $-0.30$ \\
VN100  & 1 & 0.4798 & 0.7031 & \textbf{0.4735} & & $-6.23^{*}$ & $+0.11$ \\
VN100  & 5 & 0.5441 & 0.6993 & \textbf{0.5372} & & $-3.34^{*}$ & $-1.00$ \\
S\&P500& 1 & \textbf{0.3616} & 0.7286 & 0.3654 & & $+10.39^{*}$ & $-2.98^{*}$ \\
S\&P500& 5 & \textbf{0.4226} & 0.7287 & 0.4232 & & $+1.01$      & $-3.98^{*}$ \\
\bottomrule
\end{tabular}
\end{table}

\subsection{Graph-check ablation}
Table~\ref{tab:graph} gives the snapshot QLIKE and the leave-one-out DM. Removing the GAT branch improves
QLIKE in all eight configurations (LSTM (w/o GAT) vs LSTM+GAT favors the no-graph model everywhere,
significant on VN30 and S\&P~500), so the graph does not help; at best it is QLIKE-neutral on VN100. Under
this design HAR is strongest on the Vietnamese panels, while on the large S\&P~500 the price-only LSTM (w/o
GAT) attains the best QLIKE at $h5$ (0.3582 vs HAR 0.3680).

\begin{table}[t]
\centering
\caption{Graph-check (common-date snapshot, lookback 10). QLIKE (5-seed mean) and Diebold--Mariano
LSTM+GAT-vs-\{HAR, LSTM(w/o GAT)\} on QLIKE ($*$: $p<0.05$; positive favors the second model).}
\label{tab:graph}
\begin{tabular}{llcccc c cc}
\toprule
Panel & $h$ & HAR & GARCH & w/o GAT & LSTM+GAT & & vs HAR & vs w/o-GAT \\
\midrule
VN30   & 1 & \textbf{0.3946} & 0.6500 & 0.4120 & 0.4528 & & $+9.11^{*}$ & $+6.40^{*}$ \\
VN30   & 5 & \textbf{0.4531} & 0.6420 & 0.4663 & 0.4991 & & $+6.09^{*}$ & $+4.94^{*}$ \\
VN100  & 1 & \textbf{0.4843} & 0.6210 & 0.5204 & 0.5296 & & $+4.01^{*}$ & $+0.93$ \\
VN100  & 5 & \textbf{0.5442} & 0.6157 & 0.5552 & 0.5588 & & $+1.11$     & $+0.38$ \\
S\&P500& 1 & \textbf{0.3390} & 0.3842 & 0.3401 & 0.3473 & & $+2.53^{*}$ & $+3.23^{*}$ \\
S\&P500& 5 & 0.3680 & 0.3890 & \textbf{0.3582} & 0.3701 & & $+0.59$     & $+6.87^{*}$ \\
\bottomrule
\end{tabular}
\end{table}

\subsection{Lookback variation and data scaling}
On VN30 the longer lookback (22 vs 10) does not help: LSTM (w/o GAT) QLIKE is essentially unchanged
(0.4137 vs 0.4120 at $h1$) and LSTM+GAT is slightly worse (0.4714 vs 0.4528); the graph again hurts
(DM $+8.80$ at $h1$). Reading QLIKE across panels of increasing size, under the panel-correct date-clustered
inference of Sec.~\ref{sec:haranchor}, shows a consistent data-scaling gradient: HAR is not significantly
beaten on VN30 or VN100 at any horizon, whereas on the large S\&P~500 subset the deep temporal LSTM
significantly beats HAR at every horizon (dQLIKE $+3.01\%$ at $h1$ to $+7.10\%$ at $h22$, $p<0.001$ except
$h10$ at $p=0.001$) and the convex combination beats HAR at every horizon ($p<0.001$). The significant
beat-HAR result thus emerges only once the panel is large and data-rich, and it holds under the same
inference convention under which the Vietnamese panels show no significant win.

\subsection{HAR-anchored residual and forecast-combination study (E0--E10)}\label{sec:haranchor}
A complementary question to ``does a full-target deep model beat HAR'' is whether HAR-\emph{anchored}
learning --- combining or correcting HAR with a deep expert rather than replacing it --- improves on HAR,
and whether the cross-sectional graph adds anything once HAR carries the level. We evaluate an eleven-rung
ladder: \textbf{E0} HAR; \textbf{E1}/\textbf{E2} full-target LSTM / LSTM+GAT; \textbf{E3} a frozen-expert
convex combination of HAR and the deep expert with a single mixing weight fit on the validation predictions
only; \textbf{E5}/\textbf{E6}/\textbf{E7} additive HAR-residual models whose deep branch (LSTM-only, the
graph/GAT branch, or the combined branch) predicts the HAR residual with a zero-initialised head (equal to
HAR at initialisation); \textbf{E8} a multiplicative HAR-anchored residual $(\mathrm{HAR}+\varepsilon)\cdot
\exp(\text{corr})$, positive and bounded by construction; and \textbf{E9}/\textbf{E10} static and dynamic
(regime-aware) gates scaling the residual correction. The design is the common-date snapshot design of the
graph study, extended to $h\in\{1,5,10,22\}$, with a target-overlap purge of $h$ snapshots at each boundary,
a per-horizon pooled-OLS HAR anchor fit on training rows only, and residual targets built from
expanding-window \emph{cross-fitted} HAR predictions (out-of-sample residuals, not in-sample). Five seeds;
primary metric QLIKE. Significance uses the \textbf{date-clustered} Diebold--Mariano test (one loss
differential per calendar date), the panel-correct treatment for a cross-section that shares each date.

Tables~\ref{tab:haranchor-vn30} and~\ref{tab:haranchor-vn100} report QLIKE, the out-of-sample $R^2$ relative
to HAR, and the date-clustered DM $p$-value versus HAR. \textbf{No model significantly beats HAR on VN30 or
VN100 at any horizon}: the convex combination (E3), the closest competitor, has date-clustered $p\in$
0.14--0.51 (VN30) and 0.08--0.29 (VN100), never below 0.05; the residual and gated variants have larger
$p$-values still. Point-estimate QLIKE gaps favour the hybrids at longer horizons (e.g.\ E3 reaches
$+5.85\%$ QLIKE at VN100 $h22$, 0.5816 vs HAR 0.6177), but they are within noise given the short common-date
test windows (of order 130 dates on VN30, 50 on VN100). Full-target deep models (E1/E2) underperform HAR at
short horizons ($h1$: VN30 $p<0.001$; VN100 $p=0.023$/$0.035$), whereas the HAR-anchored variants track HAR
closely at short horizons and never trail far, i.e.\ anchoring removes the short-horizon penalty.

\textbf{Large-panel result (S\&P~500).} On the data-rich S\&P~500 subset (457 nodes, $\sim$300 test dates;
Table~\ref{tab:haranchor-sp500}) the same evaluation gives the opposite verdict for the deep temporal model:
the full-target LSTM (E1) significantly beats HAR at every horizon (dQLIKE $+3.01\%$/$+4.09\%$/$+4.18\%$/%
$+7.10\%$ at $h1$/$h5$/$h10$/$h22$; $p<0.001$ except $h10$ at $p=0.001$), and the convex combination (E3)
beats HAR at every horizon ($+3.16\%$ to $+5.39\%$, $p<0.001$). The graph-carrying full model (E2) is weaker
than E1 at every horizon and fails to beat HAR at $h1$ and $h10$. This beat-HAR result holds under the same
date-clustered inference under which the Vietnamese panels show no significant win, so the difference is one
of data scale, not of inference convention; on this panel the LSTM also leads on RMSE and MAE, not only
QLIKE. The additive/graph residual rungs (E5--E8) are numerically fragile on this large cross-section
(e.g.\ E5 QLIKE 2.9978 at $h22$, E7 QLIKE 5.7090 at $h10$, vs HAR 0.4596/0.4303), so the clean beat-HAR
models on the large panel are E1 and E3.

\begin{table}[t]
\centering
\caption{HAR-anchored ladder on VN30 (common-date snapshot; 5-seed means). QLIKE (lower better),
out-of-sample $R^2$ vs HAR, date-clustered DM $p$ vs HAR. n/a: the static gate has no per-observation
differential from HAR.}
\label{tab:haranchor-vn30}
\small
\begin{tabular}{llccc}
\toprule
$h$ & Model & QLIKE & $R^2_{OOS}$ vs HAR & DM $p$ (date-clust.) \\
\midrule
1  & E0 HAR                 & 0.3946 & 0.0000  & --- \\
1  & E1 LSTM (full)         & 0.4383 & $-0.1313$ & $<0.001$ \\
1  & E2 LSTM+GAT (full)     & 0.4915 & $-0.2438$ & $<0.001$ \\
1  & E3 convex combination  & 0.3984 & $-0.0017$ & 0.237 \\
1  & E5 add.\ resid.\ (LSTM)   & 0.3946 & $-0.0001$ & 0.875 \\
1  & E6 add.\ resid.\ (graph)  & 0.3946 & $-0.0001$ & 0.882 \\
1  & E7 add.\ resid.\ (both)   & 0.3946 & 0.0000  & 0.795 \\
1  & E8 mult.\ HAR-anchored    & 0.3947 & $-0.0002$ & 0.894 \\
1  & E9 static gate         & 0.3946 & 0.0000  & n/a \\
1  & E10 dynamic gate       & 0.3946 & 0.0000  & 0.754 \\
\midrule
5  & E0 HAR                 & 0.4531 & 0.0000  & --- \\
5  & E1 LSTM (full)         & 0.4810 & $-0.0803$ & 0.004 \\
5  & E2 LSTM+GAT (full)     & 0.5185 & $-0.1603$ & $<0.001$ \\
5  & E3 convex combination  & 0.4639 & $-0.0171$ & 0.135 \\
5  & E5 add.\ resid.\ (LSTM)   & 0.4531 & $-0.0001$ & 0.994 \\
5  & E6 add.\ resid.\ (graph)  & 0.4531 & $-0.0001$ & 0.926 \\
5  & E7 add.\ resid.\ (both)   & 0.4532 & $-0.0002$ & 0.906 \\
5  & E8 mult.\ HAR-anchored    & 0.4531 & $-0.0000$ & 0.954 \\
5  & E9 static gate         & 0.4531 & 0.0000  & n/a \\
5  & E10 dynamic gate       & 0.4531 & $-0.0000$ & 0.932 \\
\midrule
10 & E0 HAR                 & 0.4849 & 0.0000  & --- \\
10 & E1 LSTM (full)         & 0.5291 & $-0.1147$ & 0.110 \\
10 & E2 LSTM+GAT (full)     & 0.5394 & $-0.1221$ & 0.002 \\
10 & E3 convex combination  & 0.4958 & $-0.0244$ & 0.322 \\
10 & E5 add.\ resid.\ (LSTM)   & 0.4853 & $-0.0002$ & 0.782 \\
10 & E6 add.\ resid.\ (graph)  & 0.4854 & $-0.0003$ & 0.560 \\
10 & E7 add.\ resid.\ (both)   & 0.4854 & $-0.0003$ & 0.588 \\
10 & E8 mult.\ HAR-anchored    & 0.4850 & $-0.0001$ & 0.851 \\
10 & E9 static gate         & 0.4849 & 0.0000  & n/a \\
10 & E10 dynamic gate       & 0.4850 & $-0.0001$ & 0.600 \\
\midrule
22 & E0 HAR                 & 0.5025 & 0.0000  & --- \\
22 & E1 LSTM (full)         & 0.5185 & $-0.0609$ & 0.527 \\
22 & E2 LSTM+GAT (full)     & 0.5603 & $-0.1141$ & 0.073 \\
22 & E3 convex combination  & 0.5137 & $-0.0416$ & 0.511 \\
22 & E5 add.\ resid.\ (LSTM)   & 0.5026 & $-0.0005$ & 0.913 \\
22 & E6 add.\ resid.\ (graph)  & 0.5026 & $-0.0005$ & 0.895 \\
22 & E7 add.\ resid.\ (both)   & 0.5027 & $-0.0006$ & 0.877 \\
22 & E8 mult.\ HAR-anchored    & 0.5027 & $-0.0010$ & 0.816 \\
22 & E9 static gate         & 0.5025 & 0.0000  & n/a \\
22 & E10 dynamic gate       & 0.5025 & $-0.0001$ & 0.896 \\
\bottomrule
\end{tabular}
\end{table}

\begin{table}[t]
\centering
\caption{HAR-anchored ladder on VN100 (common-date snapshot; 5-seed means). Columns as in
Table~\ref{tab:haranchor-vn30}.}
\label{tab:haranchor-vn100}
\small
\begin{tabular}{llccc}
\toprule
$h$ & Model & QLIKE & $R^2_{OOS}$ vs HAR & DM $p$ (date-clust.) \\
\midrule
1  & E0 HAR                 & 0.4844 & 0.0000  & --- \\
1  & E1 LSTM (full)         & 0.5371 & $-0.1056$ & 0.023 \\
1  & E2 LSTM+GAT (full)     & 0.5207 & $-0.0599$ & 0.035 \\
1  & E3 convex combination  & 0.4821 & 0.0027  & 0.242 \\
1  & E5 add.\ resid.\ (LSTM)   & 0.4842 & 0.0001  & 0.782 \\
1  & E6 add.\ resid.\ (graph)  & 0.4842 & 0.0002  & 0.603 \\
1  & E7 add.\ resid.\ (both)   & 0.4842 & 0.0002  & 0.611 \\
1  & E8 mult.\ HAR-anchored    & 0.4879 & $-0.0058$ & 0.582 \\
1  & E9 static gate         & 0.4844 & 0.0000  & n/a \\
1  & E10 dynamic gate       & 0.4843 & 0.0000  & 0.601 \\
\midrule
5  & E0 HAR                 & 0.5441 & 0.0000  & --- \\
5  & E1 LSTM (full)         & 0.5614 & $-0.0153$ & 0.244 \\
5  & E2 LSTM+GAT (full)     & 0.5582 & 0.0027  & 0.415 \\
5  & E3 convex combination  & 0.5339 & 0.0208  & 0.294 \\
5  & E5 add.\ resid.\ (LSTM)   & 0.5439 & 0.0002  & 0.746 \\
5  & E6 add.\ resid.\ (graph)  & 0.5438 & 0.0003  & 0.634 \\
5  & E7 add.\ resid.\ (both)   & 0.5438 & 0.0003  & 0.631 \\
5  & E8 mult.\ HAR-anchored    & 0.5480 & 0.0072  & 0.823 \\
5  & E9 static gate         & 0.5441 & 0.0000  & n/a \\
5  & E10 dynamic gate       & 0.5441 & 0.0001  & 0.624 \\
\midrule
10 & E0 HAR                 & 0.5985 & 0.0000  & --- \\
10 & E1 LSTM (full)         & 0.5750 & 0.0433  & 0.283 \\
10 & E2 LSTM+GAT (full)     & 0.5803 & 0.0428  & 0.455 \\
10 & E3 convex combination  & 0.5701 & 0.0441  & 0.202 \\
10 & E5 add.\ resid.\ (LSTM)   & 0.5979 & 0.0003  & 0.673 \\
10 & E6 add.\ resid.\ (graph)  & 0.5883 & 0.0119  & 0.458 \\
10 & E7 add.\ resid.\ (both)   & 0.5929 & 0.0062  & 0.459 \\
10 & E8 mult.\ HAR-anchored    & 0.6008 & 0.0160  & 0.921 \\
10 & E9 static gate         & 0.5884 & 0.0118  & 0.488 \\
10 & E10 dynamic gate       & 0.5966 & 0.0020  & 0.440 \\
\midrule
22 & E0 HAR                 & 0.6177 & 0.0000  & --- \\
22 & E1 LSTM (full)         & 0.5888 & 0.0556  & 0.164 \\
22 & E2 LSTM+GAT (full)     & 0.6031 & 0.0523  & 0.540 \\
22 & E3 convex combination  & 0.5816 & 0.0512  & 0.078 \\
22 & E5 add.\ resid.\ (LSTM)   & 0.6171 & 0.0005  & 0.765 \\
22 & E6 add.\ resid.\ (graph)  & 0.5925 & 0.0374  & 0.481 \\
22 & E7 add.\ resid.\ (both)   & 0.5922 & 0.0404  & 0.555 \\
22 & E8 mult.\ HAR-anchored    & 0.6187 & 0.0065  & 0.973 \\
22 & E9 static gate         & 0.5912 & 0.0430  & 0.569 \\
22 & E10 dynamic gate       & 0.6043 & 0.0180  & 0.480 \\
\bottomrule
\end{tabular}
\end{table}

\begin{table}[t]
\centering
\caption{HAR-anchored ladder on the S\&P~500 subset (457 nodes; 5-seed means, date-clustered DM $p$ vs HAR).
Non-fragile rungs shown; the additive/graph residual rungs E5--E8 are numerically fragile on this
cross-section and are discussed in the text.}
\label{tab:haranchor-sp500}
\small
\begin{tabular}{llccc}
\toprule
$h$ & Model & QLIKE & dQLIKE\% vs HAR & DM $p$ (date-clust.) \\
\midrule
1  & E0 HAR                 & 0.3776 & ---     & --- \\
1  & GARCH                  & 0.4843 & $-28.27$ & $<0.001$ \\
1  & E1 LSTM (full)         & 0.3662 & $+3.01$  & $<0.001$ \\
1  & E2 LSTM+GAT (full)     & 0.3826 & $-1.32$  & 0.305 \\
1  & E3 convex combination  & 0.3656 & $+3.16$  & $<0.001$ \\
\midrule
5  & E0 HAR                 & 0.4158 & ---     & --- \\
5  & GARCH                  & 0.4830 & $-16.16$ & $<0.001$ \\
5  & E1 LSTM (full)         & 0.3988 & $+4.09$  & $<0.001$ \\
5  & E2 LSTM+GAT (full)     & 0.4078 & $+1.93$  & 0.005 \\
5  & E3 convex combination  & 0.4009 & $+3.58$  & $<0.001$ \\
\midrule
10 & E0 HAR                 & 0.4303 & ---     & --- \\
10 & GARCH                  & 0.4842 & $-12.53$ & $<0.001$ \\
10 & E1 LSTM (full)         & 0.4123 & $+4.18$  & 0.001 \\
10 & E2 LSTM+GAT (full)     & 0.4230 & $+1.69$  & 0.115 \\
10 & E3 convex combination  & 0.4147 & $+3.62$  & $<0.001$ \\
\midrule
22 & E0 HAR                 & 0.4596 & ---     & --- \\
22 & GARCH                  & 0.4883 & $-6.23$  & 0.121 \\
22 & E1 LSTM (full)         & 0.4270 & $+7.10$  & $<0.001$ \\
22 & E2 LSTM+GAT (full)     & 0.4386 & $+4.59$  & 0.004 \\
22 & E3 convex combination  & 0.4349 & $+5.39$  & $<0.001$ \\
\bottomrule
\end{tabular}
\end{table}

\textbf{Graph attribution.} A graph contribution requires the graph/GAT residual (E6) or the combined
residual (E7) to beat the same-capacity no-graph residual (E5), not merely HAR. The paired date-clustered DM
$p$-values (E6 vs E5, E7 vs E5) are, by horizon: VN30 --- 0.878/0.983, 0.886/0.783, 0.761/0.782,
0.995/0.929; VN100 --- 0.705/0.787, 0.795/0.873, 0.438/0.419, 0.467/0.546; S\&P~500 --- 0.0003/0.169,
0.237/0.201, 0.154/0.015, 0.088/0.118. On the Vietnamese panels none approaches significance (all
$p\ge0.41$) and the sign is mixed; on the S\&P~500 the no-graph residual is favoured at every horizon,
significantly at $h1$ ($p=0.0003$) and $h10$ ($p=0.015$). In no panel and at no horizon does the graph
residual significantly beat the same-capacity no-graph residual, consistent with the leave-one-out ablation.

\textbf{Model-free graph screening.} Because the ablations compare specific architectures, we additionally
rule out cross-stock signal model-free. Appending a neighbour-derived signal to the three HAR features in a
train-only OLS and measuring the incremental test $R^2$ over HAR upper-bounds what any model consuming those
edges could extract linearly (leakage-safe; edges train-only). Across VN30, VN100 and the S\&P~500 subset at
every horizon the incremental test $R^2$ is at most about $1\%$ and near zero for the theoretically decisive
innovation and directed lead-lag screens; the only small positives are contemporaneous neighbour-level
signals that sit close to a shuffled-edge placebo. A single leakage-safe regression of the target on HAR plus
the mean Parkinson variance of a node's Top-5 neighbours at day $t$ --- bypassing the GAT entirely --- yields
incremental out-of-sample $R^2$ of essentially zero everywhere (VN30 $-0.0001$/$+0.0013$/$-0.0014$/$-0.0004$;
VN100 $-0.0007$/$+0.0091$/$-0.0010$/$+0.0007$). Diagnostics confirm a genuine null: the graphical-lasso edge
set barely persists train$\to$test (Top-5 Jaccard 0.09--0.17), the GAT branch is alive (a $\sim$15\%-of-HAR
correction at VN100 $h22$) yet does not beat HAR, its attention collapses to a near-uniform neighbour average
(entropy $\approx0.996$), and the graph residual shows no train-vs-test overfitting gap.

\textbf{Numerical fragility of the residual rungs.} On the large S\&P~500 subset the additive and graph
residual rungs (E5--E8) can drive predictions toward the QLIKE floor and inflate QLIKE far above HAR
(e.g.\ E5 QLIKE 2.9978 at $h22$, E7 QLIKE 5.7090 at $h10$, vs HAR 0.4596/0.4303). These fragile rungs are
reported for completeness and are not the basis of the S\&P~500 beat-HAR claim, which rests on the
well-behaved full-target LSTM (E1) and convex combination (E3) of Table~\ref{tab:haranchor-sp500}.

\textbf{Methodological note (panel-correct inference).} For a cross-sectionally dependent volatility panel,
forecast-comparison significance must use date-clustered (or panel-robust) inference --- one loss
differential per calendar date --- not a naive per-observation Diebold--Mariano over every (ticker, date)
row. All tickers share each date, so the per-observation test treats the effective sample as the number of
tickers times the number of dates, understates the loss-differential variance, and inflates the statistic by
a factor of order $\sqrt{N}$ (roughly six on VN30, ten on VN100); a QLIKE gap that is a coin-flip once
collapsed to one value per date can then read as highly significant. This panel-correct convention is used
throughout the present study.

\textbf{Relation to the per-observation tables.} The per-observation Diebold--Mariano statistics in the main
study (Table~\ref{tab:main}) flag short-horizon LSTM-vs-HAR gaps on the Vietnamese panels as significant;
that test overstates significance on this cross-sectionally dependent panel, and under the date-clustered
inference used throughout, the Vietnamese gaps of that magnitude are not significant. The point estimates are
unchanged; only their per-observation significance is superseded by the date-clustered verdict, symmetrically
for the cells where HAR nominally beats the deep model. The one beat-HAR result that survives date-clustered
inference is on the large S\&P~500 subset (Table~\ref{tab:haranchor-sp500}), not on the Vietnamese panels.

\section{Discussion}\label{sec:disc}
\textbf{HAR is hard to beat on small panels; data scale unlocks a deep advantage.} Under panel-correct
date-clustered inference no model significantly beats HAR on VN30 or VN100 at any horizon, while the same
architectures, evaluated identically, significantly beat HAR at every horizon on the large S\&P~500 subset
(the LSTM improves QLIKE by up to $+7.1\%$ at $h22$; the convex combination beats HAR at every horizon).
Because the inference convention is held fixed across panels, the difference is one of data scale: a flexible
model needs a large, data-rich cross-section to out-of-sample beat a parsimonious baseline, and HAR's edge on
the small Vietnamese panels is a small-sample phenomenon rather than a fundamental ceiling.
\textbf{The graph adds no out-of-sample value, and this is a genuine null.} No graph-carrying model beats its
same-capacity no-graph counterpart on any panel (leave-one-out ablation, paired residual attribution, and the
S\&P~500 case where the no-graph residual is significantly favoured); the model-free screen shows a
leakage-safe linear neighbour-signal regressor adds $\approx0$ incremental out-of-sample $R^2$, and
diagnostics rule out a bug or overfitting (edge set does not persist train$\to$test, attention collapses to a
uniform average, no train-vs-test overfitting gap). The cross-sectional dependence the graph encodes is
dominated by a common market factor already captured by each ticker's own HAR persistence.
\textbf{Data design and inference convention both shape the apparent verdict.} Two secondary observations
support this reading: the deep-vs-HAR point estimates depend on the data design (the same LSTM's estimates
favour it under the richer per-observation design but not under the common-date snapshot), and, more
importantly, the significance verdict depends on the inference convention (a naive per-observation test
inflates the statistic by $\sim\sqrt{N}$). Both caution against reading a beat-HAR claim from a favourable
design plus a naive test. \textbf{Metrics.} We report all five metrics rather than the one that flatters a
model; on the S\&P~500 the deep temporal model leads on QLIKE, RMSE and MAE simultaneously, while on the
Vietnamese panels the best metric differs by model (HAR on QLIKE and squared error, the multiplicative
HAR-anchored residual on MAE).

\section{Limitations}
The graph-defined snapshot design retains on the order of 130 test dates on VN30 and 50 on VN100; under
date-clustered inference these are the effective sample sizes, so even sizeable long-horizon point-estimate
gaps (e.g.\ the convex combination's $+5.85\%$ at VN100 $h22$) are not significant --- a ``not demonstrated''
rather than ``no effect at any sample size'' result. Model selection uses validation QLIKE for the
HAR-anchored study and validation MSE for the main-study point-estimate tables; we report all five metrics
rather than cherry-pick. The graph is only evaluable on the data-poorer snapshot design, though the negative
result rests on within-design comparisons plus the design-independent model-free screen; the graph consumed
here is the binary support of the graphical-lasso graph, and while the screen additionally tests signed,
weighted, innovation and directed lead-lag signals, volume-shock and sector graphs remain follow-up work. GAT
attention is quadratic in nodes (small batch needed at $\sim$500 nodes). The S\&P~500 panel is the
long-history subset of current constituents (survivorship bias); we claim the data-scaling direction and the
beat-HAR-under-scale result on this subset, not the full index or the absolute level. Vietnam is the case
study and the S\&P~500 subset the cross-market check; no universal generality is claimed, and a masked-panel
robustness check that relaxes the common-date intersection is ongoing and not included here. All numbers are
read from stored results.

\section{Conclusion}
Evaluated with panel-correct date-clustered Diebold--Mariano inference on three equity panels, HAR is very
hard to beat on the small Vietnamese markets: no model --- LSTM, LSTM+GAT, convex combination or HAR-anchored
residual --- significantly beats HAR on VN30 or VN100 at any horizon. On the large, data-rich S\&P~500 subset
the same evaluation shows the deep temporal LSTM (no graph) significantly beating HAR at every horizon, by up
to $+7.1\%$ QLIKE at $h22$ ($p<0.001$), with the convex HAR--deep combination beating HAR at every horizon;
every learned model beats GARCH. Deep-versus-HAR competitiveness therefore grows with panel size. Adding a
Graph Attention Network over a graphical-lasso graph contributes no out-of-sample value on any panel --- a
genuine null established by leave-one-out and paired-residual ablations and by a model-free neighbour-signal
screen, with diagnostics ruling out a bug or overfitting. We recommend a graph-free HAR or HAR-anchored deep
model: HAR alone on small panels, where it is not significantly beaten, and a deep temporal or
convex-combination model on large, data-rich panels, where it significantly improves on HAR.

\begin{figure}[t]
\centering
% figure generated by docs/paper/diagrams/generate_arch.py -> diagrams/soict_harlstmgat.{pdf,png,svg}
\includegraphics[width=\textwidth]{soict_harlstmgat.pdf}
\caption{LSTM+GAT: a per-node LSTM temporal branch and a GAT spatial branch over a graphical-lasso
partial-correlation graph are concatenated into an MLP head; the leave-one-out ablation removes the GAT
branch to give the main LSTM.}
\label{fig:arch}
\end{figure}

\begin{thebibliography}{99}
\bibitem{corsi2009} Corsi, F.: A simple approximate long-memory model of realized volatility. J. Financial Econometrics \textbf{7}(2), 174--196 (2009)
\bibitem{parkinson1980} Parkinson, M.: The extreme value method for estimating the variance of the rate of return. J. Business \textbf{53}(1), 61--65 (1980)
\bibitem{lstm1997} Hochreiter, S., Schmidhuber, J.: Long short-term memory. Neural Computation \textbf{9}(8), 1735--1780 (1997)
\bibitem{gat2018} Veli\v{c}kovi\'c, P., et al.: Graph attention networks. In: ICLR (2018)
\bibitem{glasso2008} Friedman, J., Hastie, T., Tibshirani, R.: Sparse inverse covariance estimation with the graphical lasso. Biostatistics \textbf{9}(3), 432--441 (2008)
\bibitem{dm1995} Diebold, F.X., Mariano, R.S.: Comparing predictive accuracy. J. Business \& Economic Statistics \textbf{13}(3), 253--263 (1995)
\bibitem{hln1997} Harvey, D., Leybourne, S., Newbold, P.: Testing the equality of prediction mean squared errors. Int. J. Forecasting \textbf{13}(2), 281--291 (1997)
\bibitem{patton2011} Patton, A.J.: Volatility forecast comparison using imperfect volatility proxies. J. Econometrics \textbf{160}(1), 246--256 (2011)
\bibitem{harq2016} Bollerslev, T., Patton, A.J., Quaedvlieg, R.: Exploiting the errors: a simple approach for improved volatility forecasting. J. Econometrics \textbf{192}(1), 1--18 (2016)
\bibitem{chen2022} Chen, Q., Robert, C.-Y.: Multivariate realized volatility forecasting with graph neural network. In: ACM ICAIF (2022). arXiv:2112.09015
\bibitem{gnnhar2025} Zhang, C., Pu, X., Cucuringu, M., Dong, X.: Forecasting realized volatility with spillover effects: perspectives from graph neural networks (2025). arXiv:2308.01419
\bibitem{audrino2025} Audrino, F., Chassot, J.: HARd to beat: the overlooked impact of rolling windows in the era of machine learning. Int. J. Forecasting (2025). arXiv:2406.08041
\end{thebibliography}

\end{document}
